% TOP_PROBLEM: {"id": 30007339, "title": "Classification of every bounded stable Allen–Cahn solution in dimension seven", "category_id": 9, "category": {"id": 9, "name": "pde", "display_name": "Partial Differential Equations", "description": "PDEs and their applications in physics and geometry.", "slug": "pde", "order_index": 9, "created_at": "2026-07-31T15:26:25.671Z"}, "status": "open"}
% FIRST_POSTED: 2026-10-06
% ATTEMPT_STATUS: solved
% OPENAI_CLAIM: {"schema_version":1,"source_repo":"https://github.com/openai/math","source_commit":"adc7f1241b42e322a6451854ab7e4b4c146bf78a","release_date":"2026-10-06","match":"full","resolution":"proved","independently_reviewed":false,"second_pass":"agreed","families":[{"family":"375","title":"De Giorgi's conjecture in dimension eight","subject":"Partial differential equations","manuscripts":[{"title":"A positive resolution of De Giorgi's conjecture in dimension eight","pdf_path":"preprints/De-Giorgis-conjecture-in-dimension-eight-September-26-2026/article.pdf","readme_path":"preprints/De-Giorgis-conjecture-in-dimension-eight-September-26-2026/README.md","date":"2026-09-26","theorem_ref":"thm:monotone (sections/00-introduction.tex:59); thm:stable (sections/00-introduction.tex:103)","statement_sources":[{"path":"preprints/De-Giorgis-conjecture-in-dimension-eight-September-26-2026/build/sections/00-introduction.tex","line":59,"label":"thm:monotone"},{"path":"preprints/De-Giorgis-conjecture-in-dimension-eight-September-26-2026/build/sections/00-introduction.tex","line":103,"label":"thm:stable"}]}],"lean":null,"reasoning_trace":null}]}
\documentclass[11pt]{article}
\usepackage[margin=25mm]{geometry}
\usepackage{fontspec}
\setmainfont{Latin Modern Roman}
\usepackage{amsmath,amssymb,mathtools}
\usepackage{unicode-math}
\setmathfont{Latin Modern Math}
\usepackage{enumitem}
\usepackage{xurl,hyperref,bookmark}
\hypersetup{colorlinks=true,urlcolor=blue}
\setcounter{secnumdepth}{0}
\setlength{\parindent}{0pt}
\setlength{\parskip}{5pt}
\setlength{\emergencystretch}{3em}
\begin{document}
\section*{Classification of every bounded stable Allen–Cahn solution in dimension seven}

\subsection{OpenAI's claim}
\textit{De Giorgi rigidity in dimension eight}: The paper claims that every strictly direction-monotone entire solution of the scalar Allen--Cahn equation in $\mathbb R^8$ is a planar heteroclinic. It also classifies stable solutions in $\mathbb R^7$ without an energy-growth assumption.

\subsection{Mathematical statement of the claimed result}
\paragraph{Mathematical claim 1.} The monotone theorem claims that if $u\in C^2(\mathbb R^8,(-1,1))$ solves $\Delta u=u^3-u$ and $\partial_8u>0$, then \[u(x)=\tanh((e\cdot x-c)/\sqrt2),\qquad e\in S^7,\quad e_8>0.\] No limits at infinity are assumed. Its stable seven-dimensional theorem states that every $v\in C^2(\mathbb R^7,[-1,1])$ solving the same equation and satisfying $\int(|\nabla\phi|^2+(3v^2-1)\phi^2)\ge0$ for compactly supported smooth $\phi$ is either $1$, $-1$ or a one-dimensional transition $\tanh((a\cdot x-b)/\sqrt2)$ with $a\in S^6$. No energy-growth assumption is used; lower-dimensional monotone cases follow by adding unused coordinates.

\subsection{Outline of the argument}
\textit{De Giorgi rigidity in dimension eight}: Finite-density stability first forces a planar profile. To remove that density assumption in dimension seven, penalized reach controls both interface bending and neighboring sheets, and stability bounds a fourth curvature moment. An ambient radial test and matched sheet amplitudes contradict unbounded interface density; the stable directional limits of the eight-dimensional monotone solution then provide the well limits needed for the planar conclusion.

\subsection{Scope relative to this problem}
\textit{De Giorgi rigidity in dimension eight}: The stable classification uses the scalar double-well equation in dimension seven. It assumes neither growth of the energy nor prescribed states at infinity and permits only the two constant wells or the explicitly planar transition profile.

With the hypotheses and conventions above, the claimed result addresses the whole mathematical assertion on this page.

\subsection{Formal verification}
\textit{De Giorgi rigidity in dimension eight}: The pinned catalogue lists no Lean formalization for this result.

This site has not run the Lean development. These links describe OpenAI’s selected formal statements and should be read with their stated scope.

\subsection{Status caveat}
OpenAI reports that an internal model produced the manuscripts. Their arguments have not been independently verified as proofs by this site. Any portion outside the documented formal scope may contain mistakes, and the catalogue label reports the announced mathematical claim.

\subsection{Sources}
\begin{itemize}
\item OpenAI, \emph{A positive resolution of De Giorgi's conjecture in dimension eight}, 2026-09-26; citation key \texttt{OAI:De-Giorgis-conjecture-in-dimension-eight-September-26-2026}. \href{https://github.com/openai/math/blob/adc7f1241b42e322a6451854ab7e4b4c146bf78a/preprints/De-Giorgis-conjecture-in-dimension-eight-September-26-2026/article.pdf}{PDF (GitHub)}; \href{https://github.com/openai/math/raw/adc7f1241b42e322a6451854ab7e4b4c146bf78a/preprints/De-Giorgis-conjecture-in-dimension-eight-September-26-2026/article.pdf}{PDF (direct)}; \href{https://github.com/openai/math/blob/adc7f1241b42e322a6451854ab7e4b4c146bf78a/preprints/De-Giorgis-conjecture-in-dimension-eight-September-26-2026/README.md}{Manuscript page and citation}; \href{https://github.com/openai/math/blob/main/preprints/De-Giorgis-conjecture-in-dimension-eight-September-26-2026/article.pdf}{Latest version}; Statement source: \href{https://github.com/openai/math/blob/adc7f1241b42e322a6451854ab7e4b4c146bf78a/preprints/De-Giorgis-conjecture-in-dimension-eight-September-26-2026/build/sections/00-introduction.tex\#L59}{thm:monotone}; Statement source: \href{https://github.com/openai/math/blob/adc7f1241b42e322a6451854ab7e4b4c146bf78a/preprints/De-Giorgis-conjecture-in-dimension-eight-September-26-2026/build/sections/00-introduction.tex\#L103}{thm:stable}.
\item \href{https://github.com/openai/math/blob/adc7f1241b42e322a6451854ab7e4b4c146bf78a/CONTENTS.md}{OpenAI catalogue entry}.
\end{itemize}
\end{document}
